% =============================================================
%  Chapter 7 — Limitations, Extensions, and Future Directions
% =============================================================

\chapter{Limitations, Extensions, and Future Directions}
\label{chap:limitations}

\providecommand{\Ham}{\mathcal{H}}
\providecommand{\Potential}{\mathcal{R}}
\providecommand{\PhaseState}{z}
\providecommand{\Pos}{q}
\providecommand{\Mom}{p}
\providecommand{\Mass}{\mathrm{M}}
\providecommand{\J}{\mathbb{J}}

% -----------------------------------------------------------------------
\section{What the Thesis Has Established}
\label{sec:ch7:achieved}
% -----------------------------------------------------------------------

The two stages, evaluated in Chapters~\ref{chap:grlsnam}--\ref{chap:experiments},
establish the following as solved:

\begin{itemize}
  \item A geometry-only structured learner in phase space
        (GRL-SNAM, Stage~1) that outperforms reactive planners and
        deep RL baselines under matched local sensing, with
        principled online correction through the Hamiltonian
        coefficient space.
  \item A material-aware risk-sensitive extension (Stage~2) that
        reduces cumulative risk exposure by 56.4\%, eliminates
        hard-hazard entries, and cuts oscillation $7\times$ each
        improvement tracing directly to a specific added energy
        term, objective, or update law.
  \item A formal enrichment principle
        (Theorem~\ref{thm:ch5:enrichment}) characterizing when and
        why each structural addition is a necessary consequence of
        the previous stage, not an arbitrary design choice.
\end{itemize}

What remains is equally important.
This chapter identifies the principal limitations of the current
formulation, maps the most natural next steps inside the existing
framework, and states the broader theoretical and empirical agenda
this work opens.
The organizing perspective throughout is:

\begin{quote}\itshape
We establish the first two rungs of an enrichment
ladder.
The ladder itself and the principle governing how new rungs are
added is the central contribution.
The remaining rungs are now clearly identified research targets.
\end{quote}

% -----------------------------------------------------------------------
\section{Current Limitations}
\label{sec:ch7:limits}
% -----------------------------------------------------------------------

\subsection{Oracle Material Access}
\label{subsec:ch7:oracle}

The most immediate limitation is that Stage~2 assumes oracle access
to the material risk field $\tilde{r}(\vec{c})$ and hazard SDF
$\phi(\vec{c})$ (Setting~2 of the thesis scope).
In real deployment, these must be inferred from raw sensor
data RGB-IR imagery, multispectral cameras, LiDAR reflection
intensity, or hyperspectral data.

The structural consequence of oracle access is not merely practical
(adding a perception module): it changes the nature of the learning
problem.
With oracle $\tilde{r}$, the soft-risk force $F_{\mathrm{soft}} =
-\lambda_s\nabla\tilde{r}$ is exact at every integration step.
With a learned belief $\hat\rho \approx \tilde{r}$, the force becomes:
\begin{equation}
  \hat{F}_{\mathrm{soft}}
  = -\lambda_s\nabla\hat\rho(\vec{c};\,\omega),
  \label{eq:ch7:approx_force}
\end{equation}
where $\omega$ parameterizes the perception model.
The gradient $\nabla\hat\rho$ now depends on both the pose $\vec{c}$
and the perception parameters $\omega$, adding a new gradient pathway
in the rollout sensitivity recursion \eqref{eq:ch4:dp_dtheta}.
Optimizing over $\omega$ jointly with the navigation parameters
$\theta$ is Stage~3 of the enrichment ladder.

\subsection{Simplified Observer Process}
\label{subsec:ch7:observer}

The current observer update \eqref{eq:ch4:observer} is:
\begin{equation}
  y_{t+1} = \Psi(y_t, \vec{q}_{t+1};\, Z^{(1)}, P_{t+1}),
  \label{eq:ch7:observer_current}
\end{equation}
which is essentially a geometric context refresh: re-centre the
obstacle extractor and risk patch at the new pose.
This is not a proper belief-update process.
It carries no uncertainty, maintains no memory across stages, and
does not model the agent's incomplete knowledge of the environment.

A richer observer would maintain a latent belief state $b_t$ over
the environment, updated sequentially:
\begin{equation}
  b_{t+1} = \mathcal{U}(b_t, \vec{q}_{t+1}, o_{t+1}),
  \qquad
  \vartheta_t = f_\xi(b_t,\, \vec{g}_k - \vec{q}_t),
  \label{eq:ch7:belief_observer}
\end{equation}
where $o_{t+1}$ are raw sensor observations and $\mathcal{U}$ is a
learned or Bayesian update operator.
The coefficient predictor $f_\xi$ would then condition on the full
belief state $b_t$ rather than the raw local context.
This extension is nontrivial: it requires a learned filtering process
co-trained with the Hamiltonian dynamics, and the gradient of
$\vartheta_t$ with respect to $b_t$ must propagate through
$\mathcal{U}$ as well as through $f_\xi$.

\subsection{Isotropic Geometry Metric}
\label{subsec:ch7:metric}

The goal attraction term in both stages uses an isotropic quadratic:
\begin{equation}
  D_{\mathcal{G}}(\vec{q}, \vec{g})
  = \|\vec{c} - \vec{g}\|^2,
  \label{eq:ch7:iso_metric}
\end{equation}
which treats all directions equally regardless of local geometry.
In practice, the "effective distance" to a goal through a narrow
corridor should weight the transverse direction more heavily than the
longitudinal; a workspace with anisotropic obstacle density should
correspondingly deform the attraction basin.

A learned anisotropic metric would replace the fixed quadratic with:
\begin{equation}
  D_{\mathcal{G}}(\vec{q}, \vec{g};\,\xi)
  = \tfrac{1}{2}(\vec{c} - \vec{g})^\top
    \Sigma_{\vec{q}}(\vec{q};\,\xi)^{-1}
    (\vec{c} - \vec{g}),
  \label{eq:ch7:aniso_metric}
\end{equation}
where $\Sigma_{\vec{q}}(\vec{q};\,\xi) \in \mathbb{S}^2_{++}$ is a
context-dependent positive-definite matrix predicted by $f_\xi$.
This is an instance of Riemannian metric learning in the configuration
space, and it would make the kinetic term $H_{\mathrm{kin}} =
\frac{1}{2}\Mom^\top\Mass(\vec{q})^{-1}\Mom$ more geometrically
faithful.
The added cost is that $\Sigma_{\vec{q}}$ must remain
positive-definite throughout training, requiring a constrained
parameterization (e.g., Cholesky factorization) and adding a new
gradient pathway through $\partial\Sigma_{\vec{q}}/\partial\xi$.

\subsection{Empirical Coverage Without Geometric Theory}
\label{subsec:ch7:coverage}

The curriculum coverage criterion \eqref{eq:ch4:coverage} counts
conquered start–goal pairs from a finite sample.
This is a practical surrogate for the richer notion of
\emph{geometric coverage}: how much of the configuration space the
learned policy can navigate stably.

A deeper coverage theory would characterize the set of start–goal
pairs $(\vec{q}_0, \vec{g})$ for which a given Hamiltonian
parameterization $\theta$ is stable and goal-reaching.
Formally, define the \emph{basin of competence}:
\begin{equation}
  \mathcal{B}(\theta)
  = \bigl\{\,(\vec{q}_0, \vec{g}) \in \mathcal{Q}^2
    \;\big|\;
    \mathcal{L}_{\mathrm{pass}}(\theta;\vec{q}_0,\vec{g})
    \leq \epsilon_\mathrm{pass}
    \;\wedge\;
    |\vec{q}_T - \vec{g}| \leq \epsilon_{\mathrm{goal}}
  \,\bigr\}.
  \label{eq:ch7:basin}
\end{equation}
The curriculum's goal is to maximize the volume
$\mathrm{vol}(\mathcal{B}(\theta))$ under the configuration-space
measure, but the current implementation uses a finite-sample proxy.
Connecting this to Riemannian volume estimates, persistent homology
of the visited state set, or information-geometric arguments about
the curvature of the energy landscape remains an open problem.

\subsection{Monitoring-Based Stability, Not Formal Certificates}
\label{subsec:ch7:stability}

The passivity condition $\mathcal{L}_{\mathrm{pass}} \leq \epsilon$
is a monitoring heuristic: it checks whether energy is non-increasing
along sampled rollouts.
It is not a Lyapunov certificate with a bounded basin of attraction.

A formal stability certificate would require bounding the Lipschitz
constant of $F_{\mathrm{tot}}(\PhaseState;\theta)$ over the state
space, which in turn requires bounding the derivatives of the learned
network $f_\xi$.
For smooth activations and bounded coefficient spaces, such bounds
exist in principle via spectral normalization or norm-bounding
techniques, but they have not been integrated into the current
training pipeline.

\subsection{Force Decomposition Identifiability}
\label{subsec:ch7:identifiability}

A subtler limitation concerns whether the factored decomposition
$H_\theta = H_{\mathrm{geom}} + H_{\mathrm{mat}}$ is identifiable
from trajectory observations.

Consider two parameter vectors $\theta \neq \theta'$ such that
$\nabla_\Pos H_\theta(\Pos,\Mom) = \nabla_\Pos H_{\theta'}(\Pos,\Mom)$
on all observed trajectories.
Then the two parameterizations produce identical dynamics and cannot
be distinguished by the training objective.
This is a standard non-identifiability problem for energy-based models,
but it takes on a specific form here:
the geometry-only and material terms could trade off coefficient mass
in ways that are invisible to the rollout loss.

More precisely, if the training data consists only of trajectories
in low-risk regions ($\tilde{r} \approx 0$), the risk heads
$(\lambda_s, \lambda_h)$ receive no gradient signal and remain at
zero regardless of their initialized values.
This is addressed practically by the curriculum (which includes
high-risk training pairs), but the theoretical condition for
unique recovery of $(\beta, \lambda_s, \lambda_h, \gamma)$ from
trajectory data analogous to identifiability conditions in inverse
RL—has not been established.

% -----------------------------------------------------------------------
\section{Immediate Extensions of the Current Framework}
\label{sec:ch7:extensions}
% -----------------------------------------------------------------------

Each of the following extensions is a direct application of the
enrichment principle (Theorem~\ref{thm:ch5:enrichment}): adding one
new energy term forces a new force channel, coefficient, objective
term, gradient pathway, and update law.

\subsection{Stage 3: Learned Material Belief Maps}
\label{subsec:ch7:stage3}

The most natural next extension replaces oracle $\tilde{r}$ with a
learned belief $\hat\rho$ inferred from raw RGB-IR observations
$Z^{(3)}$.
The enrichment chain for this transition is:
\begin{equation}
  H_{\mathrm{perc}}(\Pos;\hat\rho)
  \;\xrightarrow{(i)}\; \hat{F}_{\mathrm{soft}}
  \;\xrightarrow{(ii)}\; (\lambda_\rho, \omega)
  \;\xrightarrow{(iii)}\; \mathcal{L}_{\mathrm{perc}}
  \;\xrightarrow{(iv)}\; \nabla_{(\lambda_\rho,\omega)}\mathcal{L}
  \;\xrightarrow{(v)}\; \text{belief update law},
  \label{eq:ch7:stage3_chain}
\end{equation}
where $\omega$ parameterizes the perception backbone and
$\mathcal{L}_{\mathrm{perc}}$ includes both a navigation loss and
a perception calibration term.

Two design choices are non-trivial:
\begin{enumerate}[(i)]
  \item \textbf{Coupling perception training to navigation.}
        The perception backbone must be trained to produce
        $\hat\rho$ that is useful for the downstream force field,
        not just accurate in a pixel-wise sense.
        This requires backpropagating through the rollout into
        $\nabla_\omega\mathcal{L}_{\mathrm{nav}}$.

  \item \textbf{Uncertainty propagation.}
        A probabilistic perception model
        $\hat\rho \sim p(\rho | o, \omega)$ would allow the force
        to be $F_{\mathrm{soft}} = -\lambda_s\mathbb{E}[\nabla\rho]$
        with an additional CVaR term penalizing high-variance
        force estimates in uncertain terrain.
\end{enumerate}

\subsection{Richer Observer and Sequential Belief Filtering}
\label{subsec:ch7:filtering}

Replacing \eqref{eq:ch7:observer_current} with the full
belief-observer \eqref{eq:ch7:belief_observer} requires selecting
a representation for $b_t$.
Three natural choices are:
\begin{itemize}
  \item \textbf{Recurrent state (RNN/LSTM/Transformer).}
        $b_t$ is a hidden state; $\mathcal{U}$ is a learned
        transition function.
        Gradient backpropagation through $\mathcal{U}$ is standard.
  \item \textbf{Particle filter.}
        $b_t$ is a set of weighted particles over environment
        hypotheses.
        The gradient of $f_\xi(b_t,\cdot)$ with respect to
        $b_t$ must be approximated via reparameterization.
  \item \textbf{Gaussian belief state.}
        $b_t = (\mu_t, \Sigma_t)$, updated by an
        extended or unscented Kalman filter.
        This is the most analytically tractable but least expressive.
\end{itemize}
In all cases the update law at the belief timescale is a new fourth
loop in the hierarchy $\mathcal{U}$, operating at a timescale
between the segment loop (fast, local) and the curriculum loop
(slow, global).

\subsection{Learned Anisotropic Riemannian Metric}
\label{subsec:ch7:riemannian}

Replacing \eqref{eq:ch7:iso_metric} with
\eqref{eq:ch7:aniso_metric} requires:
\begin{enumerate}[(i)]
  \item A network head in $f_\xi$ that outputs a lower-triangular
        Cholesky factor $L(\vec{q};\xi)$ such that
        $\Sigma_{\vec{q}} = LL^\top$.
  \item A modified kinetic term
        $H_{\mathrm{kin}} = \frac{1}{2}\Mom^\top
        (L(\vec{q})L(\vec{q})^\top)^{-1}\Mom$
        that makes the symplectic structure position-dependent
        (a sub-Riemannian Hamiltonian).
  \item An additional regularizer
        $\mathcal{R}_\Sigma = w_\Sigma\|\Sigma_{\vec{q}} - I\|_F^2$
        preventing the metric from degenerating in poorly observed
        regions.
\end{enumerate}
This is a non-trivial generalization.
A position-dependent mass matrix $\Mass(\vec{q})$ does \emph{not}
alter the canonical symplectic structure on $T^*\mathcal{Q}$: the
Poisson bracket $\{f,g\} = (\nabla_{\Mom} f)^\top(\nabla_{\Pos}
g) - (\nabla_{\Pos} f)^\top(\nabla_{\Mom} g)$ remains unchanged.
What \emph{does} change is the $\dot{p}$ equation, which gains
additional terms involving $\partial\Mass^{-1}/\partial\Pos$:
\[
  \dot{p} = -\nabla_{\Pos} V(\vec{q})
            + \tfrac{1}{2}\Mom^\top
              \tfrac{\partial \Mass^{-1}(\vec{q})}{\partial \vec{q}}
              \Mom,
\]
making the geodesic structure of trajectories curvature-dependent.
Symplectic integrators designed for constant-mass Hamiltonians
must be adapted (e.g., using implicit mid-point or variational
integrators on the cotangent bundle) to correctly handle this
additional term without energy drift.

\subsection{Stronger Field-Smoothness Regularization}
\label{subsec:ch7:sobolev}

The current gradient-field regularization $\mathcal{R}_{\nabla H}$
(§\ref{sec:ch4:regularization}) uses clearance and multi-start
penalties as tractable surrogates.
The principled version is the Jacobian penalty:
\begin{equation}
  \mathcal{R}_{\mathrm{Jac}}(\theta)
  = w_J\,\frac{1}{T}\sum_t
    \bigl\|\nabla_\Pos F_{\mathrm{tot}}(\vec{q}_t,\Mom_t;\theta)\bigr\|_F^2.
  \label{eq:ch7:R_jac}
\end{equation}
Computing \eqref{eq:ch7:R_jac} exactly requires second-order
differentiation through the integrator at every training step,
which is $O(n^2)$ in the state dimension.
A natural approximation is a Sobolev-type regularizer computed
on a fixed evaluation grid:
\begin{equation}
  \mathcal{R}_{\mathrm{Sob}}(\theta)
  = w_S\sum_{\vec{q} \in \mathcal{G}}
    \bigl\|\nabla_\Pos F_{\mathrm{tot}}(\vec{q},\mathbf{0};\theta)\bigr\|^2,
  \label{eq:ch7:R_sob}
\end{equation}
where $\mathcal{G}$ is a fixed spatial grid over the configuration
space.
This is $O(|\mathcal{G}| \cdot n)$ at training time and does not
require differentiating through the integrator.

% -----------------------------------------------------------------------
\section{Broader Research Directions}
\label{sec:ch7:broader}
% -----------------------------------------------------------------------

\subsection{Beyond 2D Navigation: Richer Configuration Spaces}
\label{subsec:ch7:3d}

Both stages operate in 2D configuration spaces ($\vec{c} \in
\mathbb{R}^2$) with planar deformations.
The same framework applies in principle to:

\begin{itemize}
  \item \textbf{3D deformable navigation} - aerial or underwater
        robots with 6-DOF kinematics, where the morphology space
        $\mathcal{Q}_o$ is significantly richer and the barrier
        terms must account for 3D obstacle geometry.
  \item \textbf{Manipulation} - dexterous end-effector control
        where the configuration space is high-dimensional and
        contact constraints are dense; the IPC-style barrier
        terms are directly applicable.
  \item \textbf{Legged locomotion} - whole-body Hamiltonian control
        where ground contact, stability, and gait transitions
        are natural energy-level phenomena.
\end{itemize}

In each case the enrichment principle predicts the required
additions: new degrees of freedom in $\mathcal{Q}$ require new
energy terms, new barrier functions, and new sub-policies in the
decomposition \eqref{eq:ch3:ham}.

\subsection{More Expressive Factored Energies}
\label{subsec:ch7:richer_energy}

The current factorization
$H_\theta = H_{\mathrm{kin}} + H_{\mathrm{geom}} + H_{\mathrm{mat}}$
is flat: all terms are additive at the same level.
Richer factorizations are possible:

\begin{itemize}
  \item \textbf{Hierarchical energy.}
        $H_{\mathrm{global}} + H_{\mathrm{local}}$, where the
        global term governs broad goal-directedness and the
        local term handles fine-grained obstacle shaping.
        The meta-navigator $f_\xi$ could learn when to weight
        global vs.\ local terms based on progress signals.
  \item \textbf{Task-parameterized energy.}
        $H_\theta(\Pos,\Mom;\tau)$ where $\tau$ is a task
        descriptor (delivery vs.\ exploration vs.\ patrol),
        allowing a single parameterization to support multiple
        navigation objectives by conditioning the energy on
        the task token.
  \item \textbf{Compositional barriers.}
        Rather than summing barrier terms, one could compose
        them via a smooth max or log-sum-exp operator, which
        would prevent a single nearby obstacle from dominating
        the total barrier energy in dense clutter.
\end{itemize}

\subsection{Multi-Agent Coupled Hamiltonians}
\label{subsec:ch7:multiagent}

A natural extension of the Hamiltonian framework to multi-agent
settings would define a joint energy:
\begin{equation}
  H_{\mathrm{joint}}(\{\PhaseState^{(i)}\}_{i=1}^N)
  = \sum_i H_\theta^{(i)}(\PhaseState^{(i)})
    + \sum_{i \neq j} H_{\mathrm{inter}}^{(ij)}(\Pos^{(i)}, \Pos^{(j)}),
  \label{eq:ch7:multiagent_ham}
\end{equation}
where $H_{\mathrm{inter}}^{(ij)}$ encodes inter-agent coupling
(collision avoidance, formation maintenance, social distance).
The enrichment principle applies: adding $H_{\mathrm{inter}}$
induces new inter-agent force channels, new shared coefficients,
new multi-agent objectives, and new coordination update laws.

A necessary but not sufficient condition for compositional passivity
is that the interconnection be \emph{power-preserving}: the power
flowing out of agent $i$ through the coupling port must equal the
power flowing into agent $j$, i.e.\
$\dot H_{\mathrm{inter}}^{(ij)} = y^{(i)\top} u^{(ij)} + y^{(j)\top} u^{(ji)}$
with $u^{(ij)} + u^{(ji)} = 0$.
Non-negative coupling energy $H_{\mathrm{inter}} \geq 0$ alone is
insufficient: it prevents the coupling from being a source of energy
extraction under zero external input, but it does not preclude the
coupling port from absorbing energy from one agent and failing to
supply it back to the other.
A compositional passivity guarantee requires additionally that the
interconnection law is skew-symmetric (zero net power exchange) and
that each individual sub-Hamiltonian $H_\theta^{(i)}$ satisfies its
own passivity condition independently.
Establishing these conditions for the multi-agent navigation setting
is an open problem.

\subsection{Integration with High-Level Semantic Reasoning}
\label{subsec:ch7:llm}

Foundation-model planners
provide semantic goal specification and task decomposition at the
high level.
The Hamiltonian framework is naturally complementary: the LLM
produces a sequence of semantic goals or risk annotations that
are encoded as energy terms $H_{\mathrm{task}}$, while the
Hamiltonian dynamics ensures that these high-level intentions
translate into physically stable, risk-aware continuous trajectories.

The key open question is how to encode semantic goal specifications
(e.g., "avoid the construction zone, prefer the shaded path")
as structured energy terms with well-defined gradients and stable
dynamics.
Grounding language in energy-term coefficients—an analogue of
SayCan~\citep{ahn2022saycan} that operates at the force-field
level rather than the skill-affordance level is a promising
direction.

% -----------------------------------------------------------------------
\section{Open Theoretical Problems}
\label{sec:ch7:theory}
% -----------------------------------------------------------------------

Beyond the limitations enumerated in §\ref{sec:ch7:limits}, the
this work opens several theoretical questions that are tractable
research programs in their own right.

\paragraph{T1. Formal passivity-preserving learning guarantees.}
The current curriculum uses $\mathcal{L}_{\mathrm{pass}}$ as a
monitoring condition.
A stronger result would be a \emph{passivity-preserving update rule}:
a gradient step $\theta \leftarrow \theta - \eta\nabla\mathcal{L}$
such that $R(\theta') \succeq 0$ is guaranteed whenever
$R(\theta) \succeq 0$.
This would require either a constrained parameterization of $R_\theta$
(e.g., $R_\theta = L_\theta L_\theta^\top$ for a network $L_\theta$)
or a projection step after each gradient update.

\paragraph{T2. Multi-timescale convergence rates.}
Theorem-level convergence guarantees for the three-loop update
hierarchy (§\ref{sec:ch4:loops}) are currently absent.
Standard multi-timescale stochastic approximation theory provides asymptotic
convergence under Lipschitz assumptions, but the CVaR objective
and the curriculum re-weighting break the standard conditions.
Establishing convergence rates for the nested loop would require
either extending multi-timescale theory to risk-sensitive objectives
or using a Lyapunov function argument that couples all three loops.

\paragraph{T3. Basin-of-competence volume bounds.}
Define the basin of competence $\mathcal{B}(\theta)$ as in
\eqref{eq:ch7:basin}.
A useful theoretical result would bound
$\mathrm{vol}(\mathcal{B}(\theta))$ as a function of the
Hamiltonian's Lipschitz constant, the barrier sharpness, and the
number of training episodes.
Such a bound would connect empirical coverage to theoretical
guarantees on the policy's generalization region.

\paragraph{T4. Identifiability of the force decomposition.}
Let $\Theta(\mathcal{D})$ denote the set of parameter vectors
consistent with a trajectory dataset $\mathcal{D}$:
\begin{equation}
  \Theta(\mathcal{D})
  = \bigl\{\theta
    \;\big|\;
    \mathcal{L}(\theta;\mathcal{D}) \leq \epsilon
  \bigr\}.
  \label{eq:ch7:theta_compat}
\end{equation}
When is $\Theta(\mathcal{D})$ a singleton (up to symmetries)?
Sufficient conditions for unique recovery of
$(\beta, \lambda_s, \lambda_h, \gamma)$ would require:
\begin{enumerate}[(i)]
  \item Coverage: $\mathcal{D}$ includes trajectories in regions
        where each force channel is active.
  \item Excitation: the state–action pairs are persistently exciting
        in the sense of classical system identification.
  \item Separation: the functional forms of
        $H_{\mathrm{geom}}$ and $H_{\mathrm{mat}}$ are
        sufficiently different that their gradients cannot
        cancel over all of $\mathcal{D}$.
\end{enumerate}
Establishing formal identifiability conditions would both guide
dataset curation for training and provide a principled justification
for the multi-stage curriculum.

\paragraph{T5. Geometric Sobolev regularity of the learned field.}
The smooth navigation potential requires that
$\nabla_\Pos H_\theta$ be Lipschitz in $\Pos$.
A tighter requirement like Sobolev regularity
$F_{\mathrm{tot}} \in W^{1,\infty}(\mathcal{Q})$ would guarantee
that the flow map $\Phi^{\Delta t}_H$ is well-posed for all $\Delta t$
above a computable threshold.
Proving this for the learned network $f_\xi$ under standard
initialization and training conditions is an open problem that
connects to the spectral analysis of neural operators.

% -----------------------------------------------------------------------
\section{Final Perspective}
\label{sec:ch7:perspective}
% -----------------------------------------------------------------------

We have developed the first two stages of a continual
Hamiltonian learning framework: a geometry-only base learner and a
material-aware richer learner, connected by a formal enrichment
principle.
The experimental evidence confirms that the structured progression
is not cosmetic: each added layer resolves a specific class of failure
modes that the previous stage could not address.

Figure~\ref{fig:ch7:agenda} illustrates the full research program
that this work opens.

\begin{figure}[t]
\centering
\begin{tikzpicture}[
  font=\small\sffamily,
  rung/.style args={#1}{
    rectangle, rounded corners=5pt,
    draw=#1!70!black, fill=#1!9,
    line width=1.2pt,
    minimum width=12.0cm, minimum height=1.4cm,
    align=left, inner sep=8pt
  },
  num/.style args={#1}{
    circle,
    draw=#1!80!black,
    fill=#1!25,
    line width=1pt,
    minimum size=22pt,
    font=\bfseries\small\sffamily
  },
  arr/.style args={#1}{
    -{Stealth[length=6pt]}, line width=1.4pt, color=#1
  },
  dashed_arr/.style={
    -{Stealth[length=5pt]}, dashed,
    line width=1pt, color=gray!50
  },
  ann/.style={font=\scriptsize\sffamily\itshape, text=gray!60!black},
]

% ── Rung 1 ────────────────────────────────────────────────────────
\node[rung={teal}] (r1) at (0,0) {%
  \hspace{24pt}%
  \begin{minipage}{10.8cm}
    \textbf{Rung 1 --- Geometry-Only Structured Learning}
    \hfill\textit{\small Established (Ch.~\ref{chap:grlsnam})}\\[2pt]
    Phase-space policy, local sensing, stagewise Hamiltonian shaping,
    online geometric alignment.
  \end{minipage}
};
\node[num={teal}] at ($(r1.west)+(13pt,0)$) {\textcolor{teal!80!black}{1}};

% ── Rung 2 ────────────────────────────────────────────────────────
\node[rung={orange}, below=10pt of r1] (r2) {%
  \hspace{24pt}%
  \begin{minipage}{10.8cm}
    \textbf{Rung 2 --- Material-Aware Risk-Sensitive Learning}
    \hfill\textit{\small Established (Ch.~\ref{chap:material})}\\[2pt]
    Oracle material terms, CVaR objective, 3-loop curriculum,
    passivity-guided advancement.
  \end{minipage}
};
\node[num={orange}] at ($(r2.west)+(13pt,0)$) {\textcolor{orange!80!black}{2}};

% ── Rung 3 ────────────────────────────────────────────────────────
\node[rung={gray}, below=10pt of r2, draw=gray!70!black, fill=gray!12] (r3) {%
  \hspace{24pt}%
  \begin{minipage}{10.8cm}
    \textbf{Rung 3 --- Perception-Led Material Inference}
    \hfill\textit{\small Near-term (§\ref{subsec:ch7:stage3})}\\[2pt]
    Learned belief $\hat\rho$; perception co-training; uncertainty
    propagation through force field.
  \end{minipage}
};
\node[num={gray}] at ($(r3.west)+(13pt,0)$) {\textcolor{gray!70!black}{3}};

% ── Rung 4 ────────────────────────────────────────────────────────
\node[rung={gray}, below=10pt of r3, draw=gray!55!black, fill=gray!8] (r4) {%
  \hspace{24pt}%
  \begin{minipage}{10.8cm}
    \textbf{Rung 4 --- Anisotropic Metric and Belief-State Observer}
    \hfill\textit{\small Medium-term (§\ref{subsec:ch7:filtering}, \ref{subsec:ch7:riemannian})}\\[2pt]
    Learned $\Sigma_{\vec{q}}(\vec{q};\xi)$; sequential filtering;
    formal basin-of-competence volume bounds.
  \end{minipage}
};
\node[num={gray}] at ($(r4.west)+(13pt,0)$) {\textcolor{gray!65!black}{4}};

% ── Rung 5 ────────────────────────────────────────────────────────
\node[rung={gray}, below=10pt of r4, draw=gray!40!black, fill=gray!5] (r5) {%
  \hspace{24pt}%
  \begin{minipage}{10.8cm}
    \textbf{Rung 5 --- Self-Improving Learner in Unknown Environments}
    \hfill\textit{\small Long-term (§\ref{sec:ch7:broader})}\\[2pt]
    Coupled perception, planning, and dynamics; multi-agent
    extensions; formal convergence and identifiability guarantees.
  \end{minipage}
};
\node[num={gray}] at ($(r5.west)+(13pt,0)$) {\textcolor{gray!60!black}{5}};

% ── Arrows ────────────────────────────────────────────────────────
\draw[arr={teal!60!black}]   (r1.south) -- (r2.north);
\draw[arr={orange!60!black}] (r2.south) -- (r3.north);
\draw[dashed_arr]            (r3.south) -- (r4.north);
\draw[dashed_arr]            (r4.south) -- (r5.north);

% ── Braces ────────────────────────────────────────────────────────
\draw[
  decorate,
  decoration={brace, amplitude=7pt}, % removed mirror
  teal!60!black, line width=1pt
]
  ([xshift=8pt]r1.north east) -- ([xshift=8pt]r2.south east)
  node[midway, right=10pt, ann, text=teal!70!black]
    {\textbf{This thesis}};

\draw[
  decorate,
  decoration={brace, amplitude=7pt}, % removed mirror
  gray!50, line width=0.8pt
]
  ([xshift=8pt]r3.north east) -- ([xshift=8pt]r5.south east)
  node[midway, right=10pt, ann]
    {Future work};

\end{tikzpicture}
\caption{%
  \textbf{The full research program opened by the thesis.}
  Rungs~1 and~2 are established by the current work.
  Rungs~3--5 are identified as the natural continuation of the same
  enrichment ladder, each following the same principle:
  a new energy term forces a new force channel, coefficient,
  objective, gradient pathway, and update law.
  The ladder's organizing principle---not any single rung---is the
  framework's central contribution.
}
\label{fig:ch7:agenda}
\end{figure}

The central message is that the enrichment principle is
\emph{generative}: it does not merely describe the transition from
Rung~1 to Rung~2 retrospectively—it predicts the structure of
Rungs~3, 4, and 5 prospectively.
A researcher wishing to add perception co-training, a belief-state
observer, or an anisotropic metric knows exactly what new components
are required, because the enrichment chain (Theorem~\ref{thm:ch5:enrichment})
specifies them.

This work therefore contributes not only two working navigation
systems, but a \emph{structural language} for building richer ones.
The open problems of §\ref{sec:ch7:theory}: identifiability,
passivity-preserving updates, basin-of-competence bounds,
multi-timescale convergence are the mathematical questions whose
resolution would transform the current framework from an empirically
validated system into a theorem-supported theory of continual
structured learning.

% \begin{quote}\itshape
% The thesis has established the first two rungs of an enrichment
% ladder for continual Hamiltonian navigation.
% The ladder's organizing principle: each new energy term forces a
% coupled chain of force, coefficient, objective, and update law is
% the framework's central contribution.
% The full promise of continual Hamiltonian learning lies in ascending
% the remaining rungs: richer perception, richer geometry, and a
% deeper theoretical understanding of when and why structured
% phase-space dynamics learns reliably in unknown environments.
% \end{quote}